%% kapa: exam=1 round=24 year=2013 period=2 form=A subject=appraisal_law subject_ko=부동산관계법규 seq=36 number=76
%% source: Q-Net 감정평가사 기출문제 (cst00309) artlSeq=5073802 file=extracted/1008/5073802/2013년 제24회 감정평가사 1차시험 기출문제/2교시(A형).hwp
%% conversion: kordoc 4.24.51 -> tools/gen1.py (auto); review=proofread-sonnet; flags=none
\begin{question}{76}
측량ㆍ수로조사 및 지적에 관한 법령상 토지소유자가 하여야 하는 신청을 대신할 수 있는 자가 아닌 자는?
\begin{choices}
  \item 국가나 지방자치단체가 취득하는 토지인 경우: 해당 토지를 관리하는 행정기관의 장 또는 지방자치단체의 장
  \item 공공사업 등에 따라 학교용지의 지목으로 되는 토지인 경우: 해당 토지를 관리하는 지방자치단체의 장
  \item 공공사업 등에 따라 제방, 하천의 지목으로 되는 토지인 경우: 해당 사업의 시행자
  \item 「주택법」에 따른 공동주택의 부지인 경우: 「집합건물의 소유 및 관리에 관한 법률」에 따른 관리인(관리인이 없는 경우에는 공유자가 선임한 대표자) 또는 해당 사업의 시행자
  \item 「민법」제404조에 따른 채권자
\end{choices}
\end{question}
\answer{2}
